% foundations/principles.tex

\chapter{Architectural Principles}
\label{chap:architectural-principles}

The preceding chapter established the mathematical axioms underlying the
intrinsic description of railway alignments. Mathematics alone does not,
however, determine how engineering concepts, realizations, states, and
representations are organized. The approved Knowledge Kernel therefore
provides the architectural principles used throughout Alignment
Information Modelling (AIM).

Unlike mathematical axioms, these principles govern conceptual
responsibility and separation. They prevent different engineering layers
from being collapsed into a single geometric or computational model.

% ============================================================
\section{Intrinsic-first Formulation}
% ============================================================

\begin{principle}[Intrinsic-first Formulation]
\label{princ:intrinsic-first}
Whenever possible, alignment-related engineering knowledge is formulated
through quantities intrinsic to the alignment and parameterized by
arc-length.
\end{principle}

Intrinsic descriptions do not depend on a particular coordinate
reference system, visualization, file format, or software
implementation. Metric and spatial meaning may be added subsequently
without changing the identity of the intrinsic engineering concept.

% ============================================================
\section{Semantic and Metric Separation}
% ============================================================

\begin{principle}[Semantic--Metric Separation]
\label{princ:semantic-metric-separation}
Constructive engineering semantics and metric realization are distinct
conceptual layers.
\end{principle}

Engineering semantics describe what an engineering object is and which
constructive relations define it. Metric realization determines how that
engineering description is interpreted within a metric space.

% ============================================================
\section{Canonical Layer Separation}
% ============================================================

\begin{principle}[Canonical Layer Separation]
\label{princ:canonical-layer}
Intrinsic geometry, metric realization, spatial interpretation,
engineering realization, operational behaviour, and representation are
distinct architectural layers.
\end{principle}

Each layer answers a different engineering question. Their separation
preserves clear responsibility while allowing explicit operators and
relations to connect them.

% ============================================================
\section{CRS Independence}
% ============================================================

\begin{principle}[CRS Independence]
\label{princ:crs-independence}
Intrinsic alignment information is independent of any particular
coordinate reference system.
\end{principle}

A coordinate reference system participates in metric realization. A
change of realization context may change coordinates and metric
interpretation without redefining the underlying alignment identity or
constructive semantics.

% ============================================================
\section{State Distinction}
% ============================================================

\begin{principle}[State Distinction]
\label{princ:runtime-distinction}
Stored engineering information, evaluated engineering states, observed
states, and physically realized states are distinct.
\end{principle}

Evaluation may derive a state from engineering knowledge; observation
may describe a realized condition; and construction or maintenance may
change physical reality. None of these relations establishes identity
between the participating concepts.

% ============================================================
\section{Target and Realized Geometry}
% ============================================================

\begin{principle}[Target versus Realized Geometry]
\label{princ:target-vs-realization}
Target geometry and realized geometry are distinct engineering
conditions.
\end{principle}

Construction, maintenance, measurement, and physical response relate the
intended and realized conditions, but do not eliminate their conceptual
difference.

% ============================================================
\section{Representation Independence}
% ============================================================

\begin{principle}[Representation Independence]
\label{princ:representation-independence}
A representation is derived from engineering knowledge and is not the
engineering object or knowledge that it represents.
\end{principle}

CAD geometry, GIS layers, exchange formats, sampled polylines, rendered
views, and software data structures may communicate or process
engineering information. They do not acquire canonical authority merely
because they are stored or displayed.

% ============================================================
\section{Physical Interpretation}
% ============================================================

\begin{principle}[Physical Interpretation]
\label{princ:physical-interpretation}
Physical railway quantities arise through metric realization and spatial
interpretation rather than from intrinsic alignment semantics alone.
\end{principle}

Gravity, cant equilibrium, vehicle response, clearance, comfort, and
operational admissibility require a realized spatial context. They may
be derived from intrinsic geometry, but they are not contained in its
mathematical definition.

% ============================================================
\section{Consequences for Terminology}
% ============================================================

The preceding principles require several distinctions to remain explicit
throughout the thesis:
\[
\begin{aligned}
\text{alignment} &\neq \text{representation of an alignment},\\
\text{intrinsic geometry} &\neq \text{metric realization},\\
\text{metric realization} &\neq \text{physical realization},\\
\text{pose} &\neq \text{frame},\\
\text{coordinate transformation} &\neq \text{world-to-track mapping},\\
\text{target condition} &\neq \text{realized condition}.
\end{aligned}
\]

These principles establish the responsibility boundaries followed by the
remaining foundation chapters. Notation, state descriptions, operators,
and realization mechanisms are introduced separately and connected only
through explicit definitions and mappings.
